%------------------------------------------------------------------------------
% Redefine theorem environments at \begin{document} time
% (after Quarto's amsthm \newtheorem, to avoid "already defined" error)
% Use \DeclareTColorBox to match elegantbook.cls visual style.
%------------------------------------------------------------------------------
% This runs in the normal token stream after \begin{document},
% so # does not need doubling.
%
% Quarto generates: \begin{definition}[name]\protect\hypertarget{id}{}\label{id}
% We accept the optional [name], and Quarto's extra tokens go into the box body.
%
% First, undefine Quarto's amsthm theorem environments
\makeatletter
\let\theorem\relax\let\endtheorem\relax
\let\definition\relax\let\enddefinition\relax
\let\lemma\relax\let\endlemma\relax
\let\corollary\relax\let\endcorollary\relax
\let\proposition\relax\let\endproposition\relax
\let\axiom\relax\let\endaxiom\relax
\let\postulate\relax\let\endpostulate\relax
\let\solution\relax\let\endsolution\relax
\let\remark\relax\let\endremark\relax
\makeatother

% Now define with simple \DeclareTColorBox (O{} for optional name)
% Note: Quarto puts \protect\hypertarget{}{}\label{} tokens BEFORE the theorem body.
% These end up inside the tcolorbox body and still function correctly.

% theorem (thmstyle, ♥) — 定理
\DeclareTColorBox[auto counter,number within=chapter]{theorem}{ O{} }{
  common,thmstyle,
  title={定理~\thetcbcounter\IfValueT{#1}{ (#1)}},
}

% definition (defstyle, ♣) — 定义
\DeclareTColorBox[auto counter,number within=chapter]{definition}{ O{} }{
  common,defstyle,
  title={定义~\thetcbcounter\IfValueT{#1}{ (#1)}},
}

% lemma (thmstyle, ♥) — 引理
\DeclareTColorBox[auto counter,number within=chapter]{lemma}{ O{} }{
  common,thmstyle,
  title={引理~\thetcbcounter\IfValueT{#1}{ (#1)}},
}

% corollary (thmstyle, ♥) — 推论
\DeclareTColorBox[auto counter,number within=chapter]{corollary}{ O{} }{
  common,thmstyle,
  title={推论~\thetcbcounter\IfValueT{#1}{ (#1)}},
}

% proposition (prostyle, ♠) — 命题
\DeclareTColorBox[auto counter,number within=chapter]{proposition}{ O{} }{
  common,prostyle,
  title={命题~\thetcbcounter\IfValueT{#1}{ (#1)}},
}

% axiom (thmstyle, ♥) — 公理
\DeclareTColorBox[auto counter,number within=chapter]{axiom}{ O{} }{
  common,thmstyle,
  title={公理~\thetcbcounter\IfValueT{#1}{ (#1)}},
}

% postulate (thmstyle, ♥) — 公设
\DeclareTColorBox[auto counter,number within=chapter]{postulate}{ O{} }{
  common,thmstyle,
  title={公设~\thetcbcounter\IfValueT{#1}{ (#1)}},
}

% Starred (unnumbered) versions
\DeclareTColorBox{theorem*}{ O{} }{
  common,thmstyle,
  title={定理\IfValueT{#1}{ (#1)}},
}
\DeclareTColorBox{definition*}{ O{} }{
  common,defstyle,
  title={定义\IfValueT{#1}{ (#1)}},
}
\DeclareTColorBox{lemma*}{ O{} }{
  common,thmstyle,
  title={引理\IfValueT{#1}{ (#1)}},
}
\DeclareTColorBox{corollary*}{ O{} }{
  common,thmstyle,
  title={推论\IfValueT{#1}{ (#1)}},
}
\DeclareTColorBox{proposition*}{ O{} }{
  common,prostyle,
  title={命题\IfValueT{#1}{ (#1)}},
}

% Redefine proof, solution, remark environments (match elegantbook.cls)
\makeatletter
\let\proof\relax\let\endproof\relax
\makeatother

% Proof environment — match elegantbook.cls
\newenvironment{proof}{
  \par\noindent\textbf{\color{second}\proofname\;}
  \color{black!90}\cfs}{\par}

% Solution environment — match elegantbook.cls
\newenvironment{solution}{\par\noindent\textbf{\color{main}解} \citshape}{\par}

% Remark environment — match elegantbook.cls
\newenvironment{remark}{\noindent\textbf{\color{second}注}}{\par}

% Chinese translation for proofname
\renewcommand{\proofname}{证明}

% Override Quarto's hardcoded "Table of contents"
% Quarto template unconditionally sets \contentsname to "Table of contents"
% just before \tableofcontents. We hook into \tableofcontents to re-override.
\pretocmd{\tableofcontents}{\renewcommand{\contentsname}{目\ \ 录}}{}{}

% Bibliography header suppression is now handled by before-bib.tex partial
% (natbib-based), which patches \bibliography to use \bibname and clears
% running headers. The old CSLReferences redefinition is no longer needed.
